% =====================================================================
%  CHAPTER 6: 影响生长发育的因素（对应大纲第二章）
%  参考：往届笔记第六章 生长发育影响因素
%  参考PPT：第六章 生长发育影响因素（林苑，南京医科大学）
% =====================================================================
\chapter{影响生长发育的因素}
\setcounter{chapter}{6}
\chapterintro{
\textbf{【双横线/掌握】}遗传因素对生长发育的影响（遗传度及估计方法、全基因组关联研究（GWAS）、家族聚集性、种族差异、遗传模式、表观遗传调控、遗传环境互作）。环境因素（自然地理环境因素、化学性环境污染因素、环境内分泌干扰物、物理性环境污染因素）。\textbf{【单横线/熟悉】}社会决定因素（社会经济状况（SES）、医疗卫生保障、家庭环境六维度、文化和教育、现代媒体和娱乐方式）。\textbf{【单横线/熟悉】}行为生活方式因素（膳食模式与饮食行为、体育锻炼、睡眠行为、现代媒体影响）。\textbf{【单横线/熟悉】}生长发育长期变化的概念、主要表现、变化原因及其对人类的影响。
}

% =====================================================================
%  PART A: 遗传因素
% =====================================================================
\section{生长发育遗传因素}

\begin{defbox}
\textbf{遗传（heredity）：}形态结构、生理功能及身心发育等性状在亲代与子代间的相似性和连续性。父母双方各种基因的不同组合通过受精卵传递给后代，并通过基因的表达决定子代的各种遗传性状（如外貌、身高、健康状况和智力等）。每个儿童的基因决定了其生长发育的潜能，但这些潜能能否充分发挥则受环境和社会决定因素制约，同时也受环境-遗传交互作用影响。

\textbf{核心理念：}遗传因素决定生长发育的潜能（可能性）；环境因素制约或促进潜能的实现程度，决定发育的速度及可能达到的程度（现实性）。
\end{defbox}

% ----- A.1 生长发育遗传特点 -----
\subsection{生长发育的遗传特点}

\subsubsection{（一）遗传度及其估计方法}

\begin{keybox}
\textbf{遗传度（heritability）/ 遗传力：}用来估计某一性状在群体中有多大比例的变异是由遗传因素决定的指标。遗传度的取值范围为0$\sim$1，值越大表示遗传的作用越大，值越小表示环境的作用越大。

\textbf{遗传度的两种估计方法：}

\textbf{1. 双生子研究（Twin Study）：}
\begin{itemize}[leftmargin=*]
    \item \imp{核心逻辑：}比较同卵双生子（MZ，基因相似度100\%）与异卵双生子（DZ，基因相似度50\%）在同一性状上的一致性差异。
    \item 若\imp{同卵一致性高而异卵一致性低} → 遗传因素主导。
    \item 若\imp{异卵和同卵一致性相似} → 环境因素主导。
\end{itemize}

\textbf{2. 全基因组关联研究（Genome-Wide Association Study, GWAS）：}
\begin{itemize}[leftmargin=*]
    \item \imp{定义：}在\imp{全基因组范围内}识别与性状相关的遗传变异，并基于这些遗传变异估计遗传度的方法。
    \item \imp{遗传度缺失（missing heritability）：}GWAS研究估计的遗传度通常\imp{低于}双生子研究估计的遗传度。例如，近视的GWAS遗传度约18.4\%，而双生子遗传度约80\%。
    \item \imp{导致遗传度缺失的原因：}
    \begin{enumerate}[leftmargin=*]
        \item 双生子研究的遗传度估计偏高
        \item 微效基因的遗漏
        \item 表观遗传的贡献
        \item 基因-环境互作
    \end{enumerate}
\end{itemize}
\end{keybox}

\subsubsection{（二）不同性状的遗传度}

\begin{keybox}
\textbf{遗传度的性状差异：}
\begin{enumerate}[leftmargin=*]
    \item \imp{眼、耳鼻喉、皮肤和骨骼功能域}——遗传度较高（接近或高于0.6）。
    \item \imp{认知功能域}——遗传度中等（0.4$\sim$0.5）。
    \item \imp{生殖功能域}——遗传度较低（低于0.3）。
\end{enumerate}

\textbf{遗传度的年龄变化：}遗传度可随年龄发生变化。中国双生子队列研究（Liu Q等，2015）发现，体重、身高和体质指数（BMI）的遗传度均随着年龄增长逐渐增高。例如，BMI的加性遗传效应从0$\sim$2岁的0.13上升至15$\sim$17岁的0.54；身高的加性遗传效应从0$\sim$2岁的0.12上升至15$\sim$17岁的0.87。此现象提示\imp{遗传因素在个体越接近成熟的阶段表现得越充分}。
\end{keybox}

\subsubsection{（三）生长发育的家族、种族影响}

\begin{keybox}
\textbf{1. 生长发育的家族聚集性（familial aggregation）：}
指某些性状或疾病在家族中出现的频率高于一般人群，是亲子代际之间遗传最直接的体现。家族聚集性由\imp{父母双方遗传因素}和\imp{家庭成员间相似的环境因素}共同形成，决定了生长发育的特征、趋向、潜力和限度。成年身高与父母平均身高间的遗传度约为0.75。

\textbf{2. 生长发育的种族差异性：}
种族是指具有共同起源和共同遗传特征的人群。生长发育具有种族差异性，尤其体现在\imp{体型}、\imp{躯干和四肢长度比例}等方面。不同种族群体的体格发育水平差异不全是遗传因素作用的结果，还与各群体的发展起源、生活方式、饮食习惯及社会经济状况密切相关。
\end{keybox}

\subsubsection{（四）生长发育的遗传模式}

\begin{keybox}
\textbf{1. 单基因遗传（single gene inheritance）：}
是指表型性状仅由一对等位基因控制的现象，通常遵循\imp{孟德尔遗传规律}。单基因遗传一般决定\imp{质量性状}（即有或无），在群体中呈\imp{不连续分布}。例如，身高虽属多基因遗传性状，但也可能受单基因遗传控制——FGFR3基因突变导致软骨发育不全，患者平均身高仅约130cm（侏儒症）。

\textbf{2. 多基因遗传（polygenic inheritance）：}
是指表型性状由多个基因共同参与控制的现象，每对基因对性状的效应是微小的。多基因遗传性状除受累加基因作用外，还受\imp{环境因素}的影响。\imp{数量性状}一般由多基因遗传决定，呈\imp{连续性分布}，通常符合正态分布规律。身高、体重等绝大多数生长发育指标均属于多基因遗传性状。
\end{keybox}

% ----- A.2 表观遗传调控 -----
\subsection{生长发育的表观遗传调控}

\begin{keybox}
\textbf{表观遗传（epigenetics）：}不改变基因的DNA序列而调控基因表达的可遗传或可继承机制的统称。其特点为：DNA序列不变而表型改变；改变具有\imp{可遗传性}和\imp{可逆性}；\imp{不遵循孟德尔遗传定律}。表观遗传是环境因素和细胞内的遗传物质间发生交互作用的结果。

\textbf{表观遗传的五种主要类型：}
\begin{enumerate}[leftmargin=*]
    \item \imp{DNA甲基化（DNA methylation）}——最常见的表观遗传调控方式。
    \item \imp{组蛋白修饰（histone modification）}
    \item \imp{非编码RNA（non-coding RNA）}
    \item \imp{基因印迹（gene imprinting）}
    \item \imp{染色质重塑（chromatin remodeling）}
\end{enumerate}

\textbf{表观遗传与生长发育实例：}
母亲在孕期服用长效雌激素己烯雌酚（DES），将引起子代HOXA10基因的异常DNA甲基化，导致生殖道的发育编程改变，增加子代永久性生殖道异常以及阴道癌的发病风险。此例充分说明表观遗传在儿童少年正常生长发育及成年期疾病的发生发展中发挥着重要作用。
\end{keybox}

% ----- A.3 遗传环境互作 -----
\subsection{生长发育的遗传环境互作}

\begin{keybox}
\textbf{一、基因-环境相关性（gene-environment correlation, rGE）：}
指基因与环境共同作用和串联变化的过程，强调个体的生活环境很可能是遗传自选择的结果，即遗传影响个体对于生存环境的选择。三种类型：

\begin{enumerate}[leftmargin=*]
    \item \imp{被动基因-环境相关性（passive rGE）}——个体从其父母遗传的基因与个体成长环境的关联。例如，父母既提供基因又提供家庭环境。
    \item \imp{唤起基因-环境相关性（evocative rGE）}——个体的遗传特征引起环境中其他人的反应。例如，气质活跃的儿童可能获得更多的社交互动。
    \item \imp{主动基因-环境相关性（active rGE）}——个体主动选择或创造与自身遗传倾向相关的环境。例如，有运动天赋的儿童主动参与体育活动。
\end{enumerate}

\textbf{二、基因-环境交互作用（gene-environment interaction, G$\times$E）：}
指在不同遗传倾向下，相似的环境经历对个体的影响效应不同；同时，在不同的环境经历下，遗传对个体的影响效应也不同。

\textbf{实例——母乳喂养与FTO基因：}
纯母乳喂养对不同FTO基因型儿童的BMI变化轨迹及BMI峰值年龄的影响效应有所差异。5个月以上的纯母乳喂养可以降低FTO高风险基因型儿童肥胖发生风险。母乳喂养调节早期儿童发育，并降低不良遗传因素对生命后期超重/肥胖风险的影响。
\end{keybox}

% =====================================================================
%  PART B: 物质环境因素
% =====================================================================
\section{生长发育物质环境因素}

\begin{defbox}
\textbf{物质环境（physical environment）：}人类赖以生存的物质基础。阐明主要物质环境（自然地理气候条件、各类环境污染因素和生活环境）与儿童少年生长发育的内在联系与规律，对儿童少年预防疾病、增强体质、促进身心健康发育具有重要意义。
\end{defbox}

% ----- B.1 自然地理环境因素 -----
\subsection{自然地理环境因素}

\subsubsection{（一）地理气候因素}

\begin{keybox}
\textbf{地理气候因素（geographic-climate factors）：}对生长发育具有明显的影响，主要通过以下指标反映：日照时长、年均气温、气温年均差、平均地表温度、年降水量、平均相对湿度、海拔高度、大气压、平均水气压等。

中国历次全国学生体质与健康调研都证实了生长发育的\imp{南北差异}。不同地域人群体格发育水平的差异（尤其是长期趋势），不全是地理气候因素作用的结果，还与各群体的发展起源、生活方式、饮食习惯及社会经济状况密切相关。
\end{keybox}

\subsubsection{（二）季节因素}

\begin{keybox}
\textbf{季节对生长发育的影响：}
\begin{enumerate}[leftmargin=*]
    \item \imp{身高——春季增长最快}：3$\sim$5月的身高增长值为9$\sim$11月的2倍左右，新的骨化中心也出现较多。
    \item \imp{体重——秋季增长最快}：9$\sim$11月体重增长最快，全年体重总增长的2/3发生在9月至次年2月；炎热的夏季有些儿童体重增长很少。
    \item \imp{月经初潮}：同样受季节影响，中国女童的初潮多发生在2$\sim$3月和7$\sim$8月。
\end{enumerate}
\end{keybox}

% ----- B.2 环境污染因素 -----
\subsection{环境污染因素}

\subsubsection{（一）化学性环境污染因素}

\begin{defbox}
\textbf{化学性环境污染（environmental chemical pollution）：}由于自然灾害、工业生产、人类活动等使地区环境的化学组成发生变化，使得对人体有害的化学物质成分增加。
\end{defbox}

\begin{keybox}
\textbf{1. 空气污染：}
空气污染对儿童和少年的健康影响广泛，特别是\imp{肺功能}和\imp{神经系统发育}。孕期母体暴露于空气污染显著增加胎儿宫内生长受限和早产的风险，这些不良出生结局都可能影响儿童少年生长发育，并增加健康问题的风险。

空气污染的分类：
\begin{itemize}[leftmargin=*]
    \item \imp{室外污染}——颗粒物、二氧化氮、二氧化硫和臭氧等。来源：工业排放、汽车尾气等。
    \item \imp{室内污染}——环境烟草暴露、甲醛和挥发性有机物等。来源：燃料燃烧、建筑材料和家具等。
\end{itemize}

\textbf{2. 铅污染：}
铅是环境污染物中毒性最大的重金属之一。目前许多行业和产品生产仍大量使用铅，铅污染现状依然严峻，环境铅暴露依然是影响儿童青少年健康的主要化学性污染物。

\textbf{核心原则：任何铅暴露都可能对儿童的健康产生不良影响，因此不存在"安全"的血铅水平，理想情况下血铅水平应为零。}

血铅水平对儿童健康的影响（按严重程度递增）：
\begin{itemize}[leftmargin=*]
    \item $<$5 $\mu$g/dL：智力下降、学业成绩下降；注意力相关行为和问题发生率增加。
    \item $<$10 $\mu$g/dL：听力下降；减少出生后生长、女孩和男孩的青春期延迟。
    \item 10$\sim$39 $\mu$g/dL：神经传导减慢；血红蛋白降低、贫血。
    \item 40$\sim$79 $\mu$g/dL：腹痛、便秘、绞痛、厌食、呕吐。
    \item $\geq$80 $\mu$g/dL：严重的神经影响、抽搐、昏迷、瘫痪和死亡。
\end{itemize}
\end{keybox}

\begin{keybox}
\textbf{3. 环境内分泌干扰物（Environmental Endocrine Disrupters, EEDs）：}
对机体内天然激素的产生、释放、运输、代谢、消除、结合和功能产生干扰作用，从而破坏机体内环境稳定状态的维持，并对个体功能和发育造成影响的外源性环境化学物质。

\textbf{EEDs对儿童少年健康的主要影响：}
\begin{enumerate}[leftmargin=*]
    \item \imp{性早熟}——EEDs是儿童性早熟的直接病因或促进因素。
    \item \imp{生殖毒性}——生命早期敏感窗口期暴露，影响子代生殖器官发育和精子成熟。
    \item \imp{儿童神经发育毒性}——引起持久、不可逆的学习能力缺失和行为发育障碍。
    \item \imp{抑制儿童免疫系统的正常功能}。
\end{enumerate}

\textbf{分类：}模拟/干扰雌激素、干扰睾酮、干扰甲状腺素、干扰其他内分泌功能。

\textbf{人体暴露途径：}消化道、呼吸道和皮肤接触。
\end{keybox}

\subsubsection{（二）物理性环境污染因素}

\begin{keybox}
\textbf{物理性环境污染因素三类及其健康影响：}

\begin{enumerate}[leftmargin=*]
    \item \imp{噪声污染}——损害听觉、影响儿童视觉功能和神经心理行为发育。长期暴露对心血管、消化、内分泌等系统产生不同程度的功能紊乱甚至损伤器官。
    \item \imp{电磁辐射污染}——低强度、慢性射频辐射可影响神经系统发育、引发神经衰弱、影响视力。
    \item \imp{放射性污染}——微量的放射性物质暴露一般不影响人体健康，需达到一定剂量才会发生有害作用。
\end{enumerate}
\end{keybox}

% =====================================================================
%  PART C: 社会决定因素
% =====================================================================
\section{生长发育社会决定因素}

\begin{defbox}
\textbf{社会环境（social environment）：}人类生存及活动范围内的社会物质和精神条件的总和，涵盖社会经济文化体系、家庭因素及社会关系等，是人类长期生产劳动过程中创造和积累的物质-文化综合体，是一个与自然环境相对的概念。社会因素对生长发育产生多层次、多方面的综合效应，不仅对儿童少年体格发育具重要作用，同时对心理、智力和行为发育也产生深远影响。
\end{defbox}

% ----- C.1 社会经济和医疗卫生保健 -----
\subsection{社会经济和医疗卫生保健}

\subsubsection{（一）社会经济状况}

\begin{keybox}
\textbf{社会经济状况（Socioeconomic Status, SES）：}社会因素中应用最广泛的概念，可独立于自然环境因素，对儿童少年生长发育产生直接影响。

\textbf{SES作用的阶段性变化：}
\begin{itemize}[leftmargin=*]
    \item 当社会经济发展到一定阶段，基础物质生活水准（如食物的可及性和丰富性、卫生保健等）对儿童少年生长发育的影响程度\imp{逐渐削弱}。
    \item 而社会文化相关指标（如\imp{父母职业}、\imp{受教育程度}、\imp{家庭养育环境}）的重要性\imp{逐渐凸显}。
\end{itemize}

社会经济的发展状况与医疗卫生保健服务的提升与儿童的生长发育紧密相关。随着中国经济社会不断进步，医疗卫生保健水平稳步提升，儿童生长发育状况呈现持续向好的趋势。
\end{keybox}

\subsubsection{（二）医疗卫生保障}

\begin{defbox}
\textbf{医疗卫生保障：}政府通过财政或政策的支持，为保证居民获得最基本的公共卫生和医疗服务而提供的制度安排。不同国家或地区的经济发展水平直接影响着政府对医疗卫生事业的投入程度。卫生资源配置及卫生服务的可及性和公平性与疾病的发生、流行和人群健康状况密切相关。儿童健康状况常作为衡量国家或地区经济、文化和医疗卫生发展水平的重要指标。
\end{defbox}

% ----- C.2 家庭环境 -----
\subsection{家庭环境}

\begin{keybox}
\textbf{家庭环境六维度及其对生长发育的影响：}

家庭是儿童生活的重要场所，良好的家庭环境可促进儿童身心健康发育，对小年龄儿童的作用更为显著。

\begin{enumerate}[leftmargin=*]
    \item \imp{家庭经济状况}——家庭收入直接影响儿童的居住环境、膳食营养水平、社交活动和智力投资等。

    \item \imp{父母受教育程度}——父母受教育程度高是家庭社会经济状况中的重要正向因素。\imp{母亲的教育程度}对儿童影响尤为显著，与儿童营养状况、疾病预防、早期教育等方面密切相关。

    \item \imp{家庭结构}——在中国，核心家庭和大家庭约占家庭总数的3/4。一般来说，核心家庭和大家庭的儿童在体格生长和行为发育方面优于单亲家庭的儿童。父母离异对儿童心理健康也会产生负面影响，如职业地位和心理幸福指数较低。

    \item \imp{家庭氛围}——主要取决于家庭成员间的亲密程度。良好温馨的家庭氛围对儿童健全人格特质的形成具有显著的正向促进作用。相反，家庭氛围紧张（如父母经常争吵甚至家庭暴力）会阻碍儿童正常的心理健康发展，儿童容易出现情绪多变、社交退缩和问题行为。

    \item \imp{亲子关系}——是儿童最早建立的人际关系。积极的亲子关系使儿童感受到爱与被尊重，有助于儿童对自己、他人和周围环境产生积极、乐观的认知，对培养早期社会交往具有促进作用，为将来建立积极的同伴关系、师生关系乃至婚恋关系奠定基础。

    \item \imp{教养方式}——父母的情感支持有助于培养儿童健康积极的心理特征；而过度干预和保护可能导致内向、情绪不稳等问题。父母拒绝、否认或惩罚等养育方式，甚至使用暴力，容易导致子女形成残暴、缺乏同情心、反社会倾向等个性特征。
\end{enumerate}
\end{keybox}

% ----- C.3 文化和教育 -----
\subsection{文化和教育}

\begin{keybox}
\textbf{文化和教育的三层体系：}

社会环境中文化和教育并非独立因素，其与政治制度、经济状况、社会福利、法律法规等因素相互重叠渗透，共同影响儿童少年身心健康发展。

\begin{enumerate}[leftmargin=*]
    \item \imp{家庭教育}——个体人生教育的基础和起点，影响健康和社会适应能力。
    \item \imp{学校教育}——包括幼儿园、小学、中学和大学教育，是培养德智体美劳全面发展优秀人才的保障。
    \item \imp{社会教育}——家庭和学校之外，对儿童身心发展造成影响的一切因素，是学校、家庭教育的重要补充。
\end{enumerate}
\end{keybox}

% =====================================================================
%  PART D: 行为生活方式因素
% =====================================================================
\section{生长发育行为生活方式因素}

\begin{defbox}
\textbf{生活方式（life styles）：}人们在一定社会文化、经济、风俗、家庭等因素的影响下形成的一系列日常生活习惯和生活模式。

\textbf{行为（behaviors）：}生活方式与社会背景动态关系的体现，包括膳食、运动、睡眠、娱乐、消费、社会交往等。行为和生活方式是影响儿童青少年健康的重要因素。
\end{defbox}

% ----- D.1 膳食 -----
\subsection{膳食}

\subsubsection{（一）膳食模式}

\begin{keybox}
\textbf{膳食模式：}指膳食中不同食物的数量、比例、种类或组合，以及习惯性消费频率。膳食模式不仅反映当地人群的饮食习惯和生活水平，也反映民族的传统文化、国家的经济发展和地区的环境和资源等情况。

膳食模式与儿童生长发育密切相关。儿童少年健康的膳食模式不仅能满足其生长发育的营养需求，促进体格和智力发育，还会影响一生的食物偏好，对成年期疾病的预防与控制具有重要意义。

\textbf{《中国学龄儿童膳食指南（2022）》5条关键推荐：}
\begin{enumerate}[leftmargin=*]
    \item 主动参与食物选择和制作，提高营养素养。
    \item 吃好早餐，合理选择零食，培养健康饮食行为。
    \item 天天喝奶，足量饮水，不喝含糖饮料，禁止饮酒。
    \item 多户外活动，少视屏时间，每天60分钟以上中高强度身体活动。
    \item 定期监测体格发育，保持体重适宜增长。
\end{enumerate}
\end{keybox}

\subsubsection{（二）饮食行为}

\begin{keybox}
\textbf{健康饮食行为要求：}
\begin{enumerate}[leftmargin=*]
    \item \imp{按时进食一日三餐}——早中晚三餐的能量比为\imp{3:4:3}。
    \item \imp{吃好早餐}——不吃早餐或早餐营养不足，不仅会降低儿童少年上午的学习效率，影响其学习认知能力，还会导致体能降低，影响儿童少年的体格发育。
    \item \imp{不偏食、不挑食、不暴饮暴食}——避免盲目节食，合理选择零食，减少在外就餐频率。
\end{enumerate}
\end{keybox}

% ----- D.2 体育锻炼 -----
\subsection{体育锻炼}

\begin{keybox}
\textbf{体育锻炼：}运用各种体育手段，结合日光、空气、水等自然力和卫生措施，以发展身体、增进健康、增强体质、娱乐身心为目的的身体活动过程，是儿童少年主要体力活动的形式。

\textbf{体育锻炼对儿童少年生长发育的五大促进作用：}
\begin{enumerate}[leftmargin=*]
    \item \imp{促进身体发育}——新陈代谢及能量的消耗显著增加，改善消化吸收功能；促进肌肉、骨骼生长；关节韧带变得坚韧、更灵活。
    \item \imp{促进心血管健康}——心功能储备提高，有利于血管健康。
    \item \imp{增强呼吸功能}——呼吸功能增强，降低呼吸道感染的风险。
    \item \imp{促进心理行为发育}——全面提高儿童少年心理素质，促进其个性发展；调节儿童少年的情绪状态，有助于情绪稳定；显著增强其自尊心和自信心。
    \item \imp{降低疾病风险}——降低肥胖、心血管疾病等多种慢性病的发生风险。
\end{enumerate}
\end{keybox}

% ----- D.3 睡眠行为 -----
\subsection{睡眠行为}

\begin{keybox}
\textbf{睡眠的四大功能：}
\begin{enumerate}[leftmargin=*]
    \item \imp{消除疲劳}——使机体得到充分休息和恢复。
    \item \imp{机体复原}——促进组织修复和生长激素分泌。
    \item \imp{记忆的维持和巩固}——被形象地称为"脑的营养剂"。
    \item \imp{维持代谢和免疫平衡}——优质、充足的睡眠对于维持人体代谢和免疫平衡同样十分重要。
\end{enumerate}

对于处于生长发育期的儿童青少年来说，睡眠还同时具有\imp{促进体格和神经行为发育}的作用。

\textbf{健康睡眠习惯：}
\begin{itemize}[leftmargin=*]
    \item 保障\imp{充足的睡眠时间}。
    \item 养成\imp{定时就寝、定时晨起}的规律睡眠作息习惯。
    \item 从制度上保障学校早晨上课时间应\imp{因地制宜}，不同地理经度地区的学校可灵活地决定不同季节特别是冬春季早晨上学时间。
\end{itemize}
\end{keybox}

% ----- D.4 现代媒体和娱乐方式 -----
\subsection{现代媒体和娱乐方式}

\begin{keybox}
\textbf{一、电视对儿童少年的影响：}

\begin{itemize}[leftmargin=*]
    \item \imp{正面影响}：科普、历史、地理等节目有助于提升知识面、理解力、智力和创造力等。
    \item \imp{负面影响}：
    \begin{enumerate}[leftmargin=*]
        \item 暴力、犯罪等内容对心理健康造成不良影响，可能引发暴力、酗酒、吸毒等危险行为。
        \item 视屏时间过长影响与父母及家人的互动、沟通，不利于语言发育、情绪和认知发展。
        \item 减少体力运动时间，不利于儿童少年肌肉、骨骼、大脑等身体功能发育，影响心理和情感健康。
        \item 长时间静坐导致近视、肥胖发病率增高。
        \item 限制思维活动范围，影响思维能力、个性发展、社会适应和社会交往能力等发育。
    \end{enumerate}
\end{itemize}

\textbf{二、网络使用对儿童少年的影响：}

网络的双重影响：一方面提供丰富的学习资源和社交渠道，拓展视野；另一方面，过度使用、网络成瘾和接触不良内容对儿童少年的身心健康、学业成绩和社会适应构成严重威胁。

\textbf{三、玩具和读物对儿童少年的影响：}

\begin{itemize}[leftmargin=*]
    \item \imp{玩具}——简单知识的物质载体，训练儿童的感知觉、注意力、手眼协调、肌肉运动等，有助于促进儿童认知功能发育。
    \item \imp{读物}——"信息和知识"的物化载体，促进儿童智力、语言、情感、社会性等多方面发展。
\end{itemize}
\end{keybox}

\newpage
